% Short summary coordinated with the integrated thesis, 2026-09-09.
L'acquisition du vocabulaire expressif et de la grammaire à partir d'une expérience linguistique limitée, où les formes sont inégalement fréquentes, constitue une question centrale en sciences cognitives. Les modèles permettent de tester des hypothèses d'apprentissage ; leur portée développementale dépend de la comparaison des expériences et des comportements. Cette thèse associe une calibration des comparaisons enfants--modèles à quatre études computationnelles sur le vocabulaire, la syntaxe avec récupération d'exemples, la simulation du dialogue et l'apprentissage par récompenses dérivées des réactions des adultes. Avec des modèles et des entraînements distincts, ces études examinent les contributions de l'exposition, de l'accès aux exemples mémorisés et du feedback.

Le protocole \textsc{corpus-CDI} ajuste l'exposition estimée, la quantité de productions générées et le vocabulaire évalué pour comparer le développement expressif des enfants à celui de LSTM et de Transformers au niveau des caractères. Les modèles testés ne reproduisent pas conjointement les trajectoires enfantines de diversité lexicale, de validité des formes lexicales et de croissance du vocabulaire ; les mots peu fréquents restent particulièrement difficiles à produire. Une étude distincte augmente des modèles paramétriques par récupération des $k$ plus proches voisins. Cette intervention améliore leur sensibilité aux items syntaxiques dont les mots de contenu ont une faible fréquence moyenne et réduit partiellement l'écart lié à la fréquence lexicale, avec une contribution importante de l'information structurelle. Un benchmark d'interactions enfant--adulte montre ensuite que certaines propriétés des énoncés et tendances liées à l'âge peuvent être approchées, tandis que les conversations prolongées révèlent des écarts plus marqués. Les adultes simulés présentent un alignement sémantique plus élevé et une diversité sémantique plus faible que les adultes observés. Dans une étude indépendante d'apprentissage par renforcement, les récompenses produisent des effets différenciés : l'alignement structurel améliore le plus régulièrement la grammaticalité des énoncés générés, sans amélioration cohérente aux tests de paires minimales BLiMP et Zorro.

Ces bénéfices sélectifs coexistent avec des écarts persistants entre trajectoires et tâches. La thèse apporte des benchmarks de production et d'interaction, des interventions computationnelles contrôlées et un cadre reliant leurs résultats à des conclusions développementales circonscrites. Elle motive des théories distinguant exposition, mémoire et interaction, dont les effets doivent être testés par plusieurs mesures de l'usage du langage.
